\section{论文写作与科研工具箱}

\subsection{写作心态}
论文写作更像「有逻辑的论证」：证明工作有意义、方法合理、证据充分。句式往往相对固定，难点在逻辑链条与贡献叙事。恪守学术诚信；多听写作讲座、多模仿优秀论文结构。

\subsection{写作与协作平台}
\begin{itemize}
  \item Overleaf：组内默认 LaTeX 协作平台（本手册即可在 Overleaf 编辑）。
  \item 表格：TableGenerator 等在线工具。
  \item 润色：DeepL（翻译）、Grammarly（语法，注意学术表述仍需人工把关）。
\end{itemize}

\subsection{画图}
\begin{itemize}
  \item 流程图：draw.io（免费易上手）、Visio、diagrams；
  \item 精修矢量：Adobe Illustrator；
  \item PPT 作图后导出 PDF 插入 LaTeX，文字可清晰；
  \item AI 辅助科研示意图 MCP（需自行评估适用性）：\\
        \url{https://github.com/icebird1998/drawio-scientific-illustrator}
\end{itemize}
TIP 等期刊对图表风格有规范要求，组内相关材料请向导师/师兄索取最新版；写稿时统一字体、线宽、配色与编号。

\subsection{文献与实验周边工具}
\begin{itemize}
  \item 文献管理：Zotero；
  \item 训练可视化：Weights \& Biases；模型与数据集社区：Hugging Face；
  \item 笔记与日程：Notion / Obsidian；
  \item 预印本：arXiv（增加曝光，但需批判性阅读未经审稿内容）；
  \item 匿名开源：Anonymous GitHub（\url{https://anonymous.4open.science/}）。
\end{itemize}

\subsection{AI 读论文助手}
Paper Copilot（\url{https://papercopilot.com/}）等可用于快速扫文、提炼对比表；\textbf{结论与引用必须人工核对}，不可直接当作事实写入论文。

\subsection{投稿截止时间}
\begin{itemize}
  \item CCF 会议倒计时：\url{https://ccfddl.top/}
  \item AI deadlines：\url{https://github.com/abhshkdz/ai-deadlines} 及社区维护镜像
\end{itemize}
务必区分摘要截止与全文截止，并注意时区。

\subsection{优秀开源仓库的特征}
清晰说明、可视化结果、便于复现、排版有序；论文与 GitHub 互相引用（盲审阶段除外）。
